% GENERATED FILE. Edit canonical lessons, then run npm run generate:books.
% canonical-source-hash: fnv1a64:68839c91f11670aa
% canonical-lessons: JA-C306-aisatsu, JA-C306-yoroshiku, JA-C306-nengajou, JA-C306-funabin, JA-C306-morau, JA-R306-review-302-306

\chapter{Greetings by Post}
\label{ch:ja-greetings-by-post}
\hlchaptermodality{306}

\begin{chapteropening}
\textbf{By the end of this chapter:} \emph{I can send greetings and say what I received.}

Complete the last lesson of chapter 306: I can send greetings and say what I received.
\end{chapteropening}

\section[\texorpdfstring{\ja{あいさつ} (aisatsu)}{あいさつ (aisatsu)}]{\ja{あいさつ} (aisatsu) — a greeting}

\label{lesson:JA-C306-aisatsu}



\begin{quote}
Before the new one: say the Japanese for a delivery, then the Japanese for after a long time.
\end{quote}

\subsection*{You'll want to know: \ja{あいさつ}}
\textbf{\ja{あいさつ}} — \emph{aisatsu} — \textquotedblleft{}a greeting\textquotedblright{}.

\begin{quote}
\textbf{\ja{あさ、となりの} \ja{おばあさんに} \ja{あいさつしました。}} — \emph{asa, tonari no obāsan ni aisatsu shimashita} — \textquotedblleft{}In the morning I said hello to the old lady next door.\textquotedblright{}
\end{quote}

\begin{cousinweb}
\textbf{\ja{あいさつ}} (\emph{aisatsu}) is Sino-Japanese (kango): two Chinese characters read with their on'yomi, readings borrowed from Chinese. They are \emph{ai}, \textquotedblleft{}push open\textquotedblright{} (Middle Chinese \emph{'eajX}) and \emph{satsu}, \textquotedblleft{}draw close\textquotedblright{} (\emph{tsat}) (KANJIDIC2; Middle Chinese in Baxter's transcription).
\end{cousinweb}

\begin{grammarlens}[title={What you've built}]
One more word for telling what happened.
\end{grammarlens}

\subsection*{Your turn}
\begin{itemize}
\raggedright
  \item \emph{Say it:} \emph{aisatsu}
  \item \emph{Say it:} \emph{aisatsu}, once more
  \item \emph{Say it:} \emph{asa, tonari no obāsan ni aisatsu shimashita} — In the morning I said hello to the old lady next door.
  \item \emph{Recall:} say the Japanese for a delivery, then the Japanese for after a long time, then say \emph{aisatsu} again
\end{itemize}

\begin{tcolorbox}[breakable,colback=teal!4,colframe=teal!35!black,title={Before you move on}]
What does \ja{あいさつ} mean? (\textquotedblleft{}A greeting\textquotedblright{}.) Say it once more.
\end{tcolorbox}
\section[\texorpdfstring{\ja{よろしく} (yoroshiku)}{よろしく (yoroshiku)}]{\ja{よろしく} (yoroshiku) — kind regards; please remember me}

\label{lesson:JA-C306-yoroshiku}



\begin{quote}
Before the new one: say the Japanese for after a long time, then the Japanese for a greeting.
\end{quote}

\subsection*{You'll want to know: \ja{よろしく}}
\textbf{\ja{よろしく}} — \emph{yoroshiku} — \textquotedblleft{}kind regards; please remember me\textquotedblright{}.

\begin{quote}
\textbf{\ja{せんせいに} \ja{よろしく} \ja{いって} \ja{ください。}} — \emph{sensei ni yoroshiku itte kudasai} — \textquotedblleft{}Please give my regards to your teacher.\textquotedblright{}
\end{quote}

\begin{cousinweb}
\textbf{\ja{よろしく}} (\emph{yoroshiku}) is a native Japanese word (wago), inherited rather than borrowed from Chinese, and it is usually written in kana (JMdict).
\end{cousinweb}

\begin{grammarlens}[title={What you've built}]
One more word for telling what happened.
\end{grammarlens}

\subsection*{Your turn}
\begin{itemize}
\raggedright
  \item \emph{Say it:} \emph{yoroshiku}
  \item \emph{Say it:} \emph{yoroshiku}, once more
  \item \emph{Say it:} \emph{sensei ni yoroshiku itte kudasai} — Please give my regards to your teacher.
  \item \emph{Recall:} say the Japanese for after a long time, then the Japanese for a greeting, then say \emph{yoroshiku} again
\end{itemize}

\begin{tcolorbox}[breakable,colback=teal!4,colframe=teal!35!black,title={Before you move on}]
What does \ja{よろしく} mean? (\textquotedblleft{}Kind regards; please remember me\textquotedblright{}.) Say it once more.
\end{tcolorbox}
\section[\texorpdfstring{\ja{ねんがじょう} (nengajō)}{ねんがじょう (nengajō)}]{\ja{ねんがじょう} (nengajō) — a New Year's card}

\label{lesson:JA-C306-nengajou}



\begin{quote}
Before the new one: say the Japanese for a greeting, then the Japanese for kind regards.
\end{quote}

\subsection*{You'll want to know: \ja{ねんがじょう}}
\textbf{\ja{ねんがじょう}} — \emph{nengajō} — \textquotedblleft{}a New Year's card\textquotedblright{}.

\begin{quote}
\textbf{\ja{しょうがつに} \ja{ねんがじょうが} \ja{とどきました。}} — \emph{shōgatsu ni nengajō ga todokimashita} — \textquotedblleft{}Greeting cards arrived for the new year.\textquotedblright{}
\end{quote}

\begin{cousinweb}
\textbf{\ja{ねんがじょう}} (\emph{nengajō}) is Sino-Japanese (kango): three Chinese characters read with their on'yomi, readings borrowed from Chinese. They are \emph{nen}, \textquotedblleft{}year\textquotedblright{} (Middle Chinese \emph{nen}), \emph{ga}, \textquotedblleft{}congratulations\textquotedblright{} (\emph{haH}) and \emph{jō}, \textquotedblleft{}form\textquotedblright{} (\emph{dzrjangH}) (KANJIDIC2; Middle Chinese in Baxter's transcription).
\end{cousinweb}

\begin{grammarlens}[title={What you've built}]
One more word for telling what happened.
\end{grammarlens}

\subsection*{Your turn}
\begin{itemize}
\raggedright
  \item \emph{Say it:} \emph{nengajō}
  \item \emph{Say it:} \emph{nengajō}, once more
  \item \emph{Say it:} \emph{shōgatsu ni nengajō ga todokimashita} — Greeting cards arrived for the new year.
  \item \emph{Recall:} say the Japanese for a greeting, then the Japanese for kind regards, then say \emph{nengajō} again
\end{itemize}

\begin{tcolorbox}[breakable,colback=teal!4,colframe=teal!35!black,title={Before you move on}]
What does \ja{ねんがじょう} mean? (\textquotedblleft{}A New Year's card\textquotedblright{}.) Say it once more.
\end{tcolorbox}
\section[\texorpdfstring{\ja{ふなびん} (funabin)}{ふなびん (funabin)}]{\ja{ふなびん} (funabin) — surface mail; sea mail}

\label{lesson:JA-C306-funabin}



\begin{quote}
Before the new one: say the Japanese for kind regards, then the Japanese for a New Year's card.
\end{quote}

\subsection*{You'll want to know: \ja{ふなびん}}
\textbf{\ja{ふなびん}} — \emph{funabin} — \textquotedblleft{}surface mail; sea mail\textquotedblright{}.

\begin{quote}
\textbf{\ja{はこを} \ja{ふなびんで} \ja{おくりました。}} — \emph{hako o funabin de okurimashita} — \textquotedblleft{}I sent the box by sea mail.\textquotedblright{}
\end{quote}

\begin{cousinweb}
\textbf{\ja{ふなびん}} (\emph{funabin}) is a mixed word, native and Sino-Japanese: \emph{funa}, \textquotedblleft{}ship\textquotedblright{}, is native Japanese, the form \emph{fune} takes in compounds, and \emph{bin}, \textquotedblleft{}letter\textquotedblright{}, is a Chinese character read with its on'yomi, a reading borrowed from Chinese (Middle Chinese \emph{bjienH}) (JMdict, KANJIDIC2; Middle Chinese in Baxter's transcription).
\end{cousinweb}

\begin{grammarlens}[title={What you've built}]
One more word for telling what happened.
\end{grammarlens}

\subsection*{Your turn}
\begin{itemize}
\raggedright
  \item \emph{Say it:} \emph{funabin}
  \item \emph{Say it:} \emph{funabin}, once more
  \item \emph{Say it:} \emph{hako o funabin de okurimashita} — I sent the box by sea mail.
  \item \emph{Recall:} say the Japanese for kind regards, then the Japanese for a New Year's card, then say \emph{funabin} again
\end{itemize}

\begin{tcolorbox}[breakable,colback=teal!4,colframe=teal!35!black,title={Before you move on}]
What does \ja{ふなびん} mean? (\textquotedblleft{}Surface mail; sea mail\textquotedblright{}.) Say it once more.
\end{tcolorbox}
\section[\texorpdfstring{\ja{もらう} (morau)}{もらう (morau)}]{\ja{もらう} (morau) — to receive; to get}

\label{lesson:JA-C306-morau}



\begin{quote}
Before the new one: say the Japanese for a New Year's card, then the Japanese for surface mail.
\end{quote}

\subsection*{You'll want to know: \ja{もらう}}
\textbf{\ja{もらう}} — \emph{morau} — \textquotedblleft{}to receive; to get\textquotedblright{}.

\begin{quote}
\textbf{\ja{そふから} \ja{ねんがじょうを} \ja{もらいました。}} — \emph{sofu kara nengajō o moraimashita} — \textquotedblleft{}I got a greeting card from my grandfather.\textquotedblright{}
\end{quote}

Polite \textbf{\ja{もらいます}}, \emph{moraimasu}; a godan verb. Past (polite): \textbf{\ja{いもうとが} \ja{はがきを} \ja{もらいました。}}, \emph{imōto ga hagaki o moraimashita}.

\begin{cousinweb}
\textbf{\ja{もらう}} (\emph{morau}) is a native Japanese word (wago), inherited rather than borrowed from Chinese; its kanji was attached later as a kun'yomi (JMdict).
\end{cousinweb}

\begin{grammarlens}[title={What you've built}]
One more word for telling what happened.
\end{grammarlens}

\subsection*{Your turn}
\begin{itemize}
\raggedright
  \item \emph{Say it:} \emph{morau}
  \item \emph{Say it:} \emph{morau}, once more
  \item \emph{Say it:} \emph{sofu kara nengajō o moraimashita} — I got a greeting card from my grandfather.
  \item \emph{Recall:} say the Japanese for a New Year's card, then the Japanese for surface mail, then say \emph{morau} again
\end{itemize}

\begin{tcolorbox}[breakable,colback=teal!4,colframe=teal!35!black,title={Before you move on}]
What does \ja{もらう} mean? (\textquotedblleft{}To receive; to get\textquotedblright{}.) Say it once more.
\end{tcolorbox}
\section[(dialogue)]{Looking back: from Ghost Stories to Greetings by Post}

\label{lesson:JA-R306-review-302-306}



\begin{quote}
Say the Japanese for a ghost and a spook.
\end{quote}

\subsection*{Your turn}
Say these aloud:

\begin{itemize}
\raggedright
  \item a ghost, a spook, to bear a grudge against, a legend, a god
  \item a dragon, a marsh, to float up, a small wayside shrine, auspicious
  \item thanks, repaying a kindness, a trap, to peep, needy
  \item the post, an address (on a letter), express delivery, a delivery, after a long time
  \item a greeting, kind regards, a New Year's card, surface mail, to receive
\end{itemize}

\begin{tcolorbox}[breakable,colback=teal!4,colframe=teal!35!black,title={Before you move on}]
A ghost? (\textbf{\ja{ゆうれい}}, \emph{yūrei}.) A spook? (\textbf{\ja{おばけ}}, \emph{obake}.) To bear a grudge against? (\textbf{\ja{うらむ}}, \emph{uramu}.) A legend? (\textbf{\ja{でんせつ}}, \emph{densetsu}.) A god? (\textbf{\ja{かみさま}}, \emph{kamisama}.) A dragon? (\textbf{\ja{りゅう}}, \emph{ryū}.) A marsh? (\textbf{\ja{ぬま}}, \emph{numa}.) To float up? (\textbf{\ja{うかぶ}}, \emph{ukabu}.) A small wayside shrine? (\textbf{\ja{ほこら}}, \emph{hokora}.) Auspicious? (\textbf{\ja{めでたい}}, \emph{medetai}.) Thanks? (\textbf{\ja{おれい}}, \emph{orei}.) Repaying a kindness? (\textbf{\ja{おんがえし}}, \emph{ongaeshi}.) A trap? (\textbf{\ja{わな}}, \emph{wana}.) To peep? (\textbf{\ja{のぞく}}, \emph{nozoku}.) Needy? (\textbf{\ja{まずしい}}, \emph{mazushii}.) The post? (\textbf{\ja{ゆうびん}}, \emph{yūbin}.) An address (on a letter)? (\textbf{\ja{あてさき}}, \emph{atesaki}.) Express delivery? (\textbf{\ja{そくたつ}}, \emph{sokutatsu}.) A delivery? (\textbf{\ja{はいたつ}}, \emph{haitatsu}.) After a long time? (\textbf{\ja{ひさしぶり}}, \emph{hisashiburi}.) A greeting? (\textbf{\ja{あいさつ}}, \emph{aisatsu}.) Kind regards? (\textbf{\ja{よろしく}}, \emph{yoroshiku}.) A New Year's card? (\textbf{\ja{ねんがじょう}}, \emph{nengajō}.) Surface mail? (\textbf{\ja{ふなびん}}, \emph{funabin}.) To receive? (\textbf{\ja{もらう}}, \emph{morau}.)
\end{tcolorbox}
